\scanpage{19}
\margin{\unsure{$2^{\text{è}}$ fll|2 fll|2 bl|2 bis}}
\inserted{d'où $S_{(t)}\to S_{(s)}$.} On va définir un \deleted{\unsure{l'}}isom $A_{(\deleted{\unsure{t}}s)}\simeq B\times_{\deleted{\unsure{V}}S}S_{(\deleted{\unsure{t}}s)}$, tel que \deleted{que} après le $S_{(t)}\to S_{(s)}$, on retrouve l'isom $\deleted{\unsure{B}}A_{(t)}\simeq B\times_S S_{(t)}$ donné. \deleted{Un tel isom}

\uline{Lemme} $S$ schéma noeth., \inserted{\deleted{[réduit]}} $t\in S$, $B$ schéma abélien sur $S$. Alors $\exists$ voisinage ouvert $V$ de $t$, \deleted{\unsure{dont les composantes} \inserted{\deleted{\unsure{au voisin}}}} tel que pour \deleted{\unsure{que} $\forall\,s\in V$,} tout voisinage ouvert $\deleted{\unsure{W}}V'$ de $t$, $V'\subset V$, tout schéma abélien $B'$ sur $V'$, et tout isom $\varphi_{\unsure{(t)|t}}: B'\times_S\operatorname{Spec}\mathcal O_{S,t}\simeq B\times_S\operatorname{Spec}\mathcal O_{S,t}$, $\varphi_{(t)}$ se prolonge de façon unique en un isom
\[
\varphi:B'\simeq B\mid V'.
\]

L'unicité sera garantie dès que toute composante \unsure{ponctuelle|ponctuelles} de $V$ rencontre $t$, i.e. que $V$ est contenu ds la réunion des comp. irréd. de $S$ ne contenant pas $t$.

\deleted{Soit $S'\to S$} On peut supposer $S$ affine. \inserted{Supposons d'abord $S$ réduit.} Soit $S'\to S$ fini, tel que $S'\times_S S_{(s)}\to S_{(s)}$ soit birationnel, et $S'$ réduit et géom. unibr. Alors
% archive sheet number: 18

